\chapter*{B1\quad 《流体力学与流体机械》试题库（四）}
\addcontentsline{toc}{chapter}{B1\quad《流体力学与流体机械》试题库（四）}
\markboth{B1\quad 《流体力学与流体机械》试题库（四）}{}

\begin{tishi}
本库为旧名《流体力学与流体机械》的成套试题，题型与分值：
\textbf{选择题 15 题 $\times$ 2 分 $=$ 30 分}、\textbf{问答题 2 题 $\times$ 6 分 $=$ 12 分}、
\textbf{计算题 4 题共 58 分}（15 分、15 分、12 分、16 分），合计 100 分。

原卷把答案字母直接印在题干括号内（如"1、在水力学中，单位质量力是指（ c C ）"），
本册已把答案字母从题干中\textbf{去掉}，统一移到书末参考答案。

原卷最后两题（叶片泵允许吸上真空高度 $H_s$、两台风机并联风量）属\textbf{泵与风机}内容，
对应教材第 11--14 章、不在 818 初试范围，本册按五星制的最低档标注并注明"（流体机械）"。
\end{tishi}

\section*{一、选择题（15 题 $\times$ 2 分 $=$ 30 分）}

\begin{ti}

\item 在水力学中，单位质量力是指（\quad）。
\begin{choices}
\item 单位面积液体受到的质量力
\item 单位体积液体受到的质量力
\item 单位质量液体受到的质量力
\item 单位重量液体受到的质量力
\end{choices}
\xstarfull{3}\yiju{E4 2015 年试题填空第 1 题（重度与动力黏度换算，须先明确单位质量力与重度的关系）；E5 题库选择题第 1 题；E1 slide33"第一章绪论＝重点掌握基础概念"}\nandu{易}

\item 下列流体属于牛顿流体的是（\quad）。
\begin{choices}
\item 汽油
\item 纸浆
\item 血液
\item 沥青
\end{choices}
\xstarfull{3}\yiju{E4 2026 年回忆版计算题第 1 题（牛顿流体、切应力）；E5 题库选择题第 2 题；教材式 1.9、1.10（E6）}\nandu{易}

\item 流体处于平衡状态的必要条件是（\quad）。
\begin{choices}
\item 流体无黏性
\item 流体黏度大
\item 体积力有势
\item 流体正压
\end{choices}
\xstarfull{2}\yiju{E5 题库选择题第 3 题"体积力有势"；E6 教材第 2 章静压强特性；E4 历年真题未直接出现，故压至两星}\nandu{易}

\item 欧拉法研究（\quad）的变化情况。
\begin{choices}
\item 每个流体质点的运动参数
\item 每个质点的轨迹
\item 每个空间点上的流体质点的流速
\item 通过空间上质点的运动参数
\end{choices}
\xstarfull{2}\yiju{E5 题库选择题第 4 题；E1 slide33"第一章绪论＝重点掌握基础概念"；E4 历年真题未直接出现}\nandu{易}

\item 在（\quad）流动中，流线和迹线重合。
\begin{choices}
\item 无旋
\item 有旋
\item 定常
\item 非定常
\end{choices}
\xstarfull{2}\yiju{E5 题库选择题第 5 题；E2"第 1--3 章偏概念，掌握基本公式与解释即可"；E4 历年真题未直接成题}\nandu{易}

\item 运用沿总流的伯努利方程时所选取的两个断面（\quad）。
\begin{choices}
\item 可以是任何断面
\item 必须是均流断面
\item 之间可以有急变流
\item 之间必须是缓变流
\end{choices}
\xstarfull{3}\yiju{E4 2022 年 818 单选第 6 题（伯努利方程所选计算点位置）；E5 题库选择题第 6 题；教材伯努利方程适用条件（E6）}\nandu{中}

\item 一密闭容器内下部为水，水的重度为 $\gamma$；上部为空气，空气的压强为 $p_0$。容器由静止状态自由下落，在下落过程中容器内水深为 $h$ 处的压强为（\quad）。
\begin{choices}
\item $p_0+\gamma h$
\item $p_0$
\item $0$
\item 以上都不正确
\end{choices}
\xstarfull{2}\yiju{E5 题库选择题第 7 题"密闭容器自由下落 $p=p_0$"；E6 教材第 2 章欧拉平衡微分方程；E4 历年真题未直接出现}\nandu{中}

\item 虹吸管最高处的压强（\quad）。
\begin{choices}
\item 大于大气压
\item 等于大气压
\item 小于大气压
\item 无法确定
\end{choices}
\xstarfull{3}\yiju{E4 2022 年 818 单选第 7 题（虹吸管最高处压强与大气压的关系）；E5 题库选择题第 8 题；E1 slide16"选择/判断考概念辨析"}\nandu{易}

\item 有压管道的总水头线与测压管水头线的基本规律是（\quad）。
\begin{choices}
\item 总水头线是沿程下降的
\item 测压管水头线沿程可升可降
\item 测压管水头线是沿程下降的
\item 总水头线和测压管水头线沿程可升可降
\end{choices}
\xstarfull{3}\yiju{E4 2022 年 818 单选第 3 题（测压管水头线坡度小于 0 的沿程变化）；E5 题库选择题第 9 题；E6 教材第 4 章总水头线与测压管水头线}\nandu{中}

\item 圆管和正方形管道的断面面积、长度、相对粗糙度都相等，流动属紊流粗糙区，通过的流量相等，则两种形状管道沿程水头损失之比为（\quad）。
\begin{choices}
\item 1.255
\item 0.886
\item 3.125
\item 0.258
\end{choices}
\xstarfull{3}\yiju{E4 题库与 2022 年回忆版均出现"圆管与方管沿程损失比 0.886"；E5 题库选择题第 10 题（答案 0.886）；E6 教材第 5 章沿程阻力系数}\nandu{中}

\item 运动黏滞系数的量纲是（\quad）。
\begin{choices}
\item $\mathrm{L}/\mathrm{T}^2$
\item $\mathrm{L}/\mathrm{T}^3$
\item $\mathrm{L}^2/\mathrm{T}$
\item $\mathrm{L}^3/\mathrm{T}$
\end{choices}
\xstarfull{3}\yiju{E4 2015 年试题填空第 1 题（运动黏度与动力黏度的量纲换算）；E5 题库选择题第 11 题；E1 slide33"第六章＝量纲转化"}\nandu{易}

\item 在有压管流中研究阻力问题进行模型试验，应选择的相似准则是（\quad）。
\begin{choices}
\item 弗劳德准则
\item 雷诺准则
\item 欧拉准则
\item 其他准则
\end{choices}
\xstarfull{3}\yiju{E4 2015 年试题填空"雷诺准则"；E5 题库选择题第 12 题；E1 slide33"第六章相似理论与量纲分析＝量纲转化、相似理论"}\nandu{易}

\item 流体作无旋运动的特征是（\quad）。
\begin{choices}
\item 所有流线都是直线
\item 所有迹线都是直线
\item 任意流体元的角变形为零
\item 任意一点的涡量都为零
\end{choices}
\xstarfull{3}\yiju{E4 2026 年回忆版计算题第 4 题（证明势函数存在性，前提即无旋）；E5 题库选择题第 13 题；E2"第六到第八章只有少量填空/判断/选择，掌握基本概念"}\nandu{易}

\item 叶片泵的基本性能参数允许吸上真空高度 $H_s$ 反映（\quad）。\textbf{（流体机械，复试范围）}
\begin{choices}
\item 离心泵的吸水性能
\item 离心泵的吸水管路大小
\item 离心泵的进水口位置
\item 离心泵叶轮的进水口性能
\end{choices}
\xstarfull{1}\yiju{超纲：流体机械（泵与风机），赵琴第 11--14 章，复试范围；E5 题库选择题第 14 题}\nandu{超纲}

\item 两台同型号风机在外界条件相同的情况下并联工作，并联后在并联工况点上的风量比一台风机单独工作时的风量（\quad）。\textbf{（流体机械，复试范围）}
\begin{choices}
\item 成倍增加
\item 增加幅度不明显
\item 大幅度增加，但不是成倍增加
\item 不增加
\end{choices}
\xstarfull{1}\yiju{超纲：流体机械（泵与风机），赵琴第 11--14 章，复试范围；E5 题库选择题第 15 题}\nandu{超纲}

\end{ti}

\section*{二、问答题（2 题 $\times$ 6 分 $=$ 12 分）}

\begin{ti}

\item （6 分）什么是流体的流动性与黏滞性？什么是流体的连续介质模型？
\xstarfull{4}\yiju{E5 题库问答题第 1 题；E4 2022 年 818 单选第 1 题考"连续介质假设的含义"；E1 slide33"第一章绪论＝重点掌握基础概念"；E1 slide16"简答考名词定义"}\nandu{易}

\item （6 分）厂家所提供的泵或风机的性能参数是在标准条件下试验得到的。
\begin{xiaowen}
\item 对于一般风机而言，所谓标准条件是指什么条件？
\item 以角标"0"代表样本条件，$B$ 为当地大气压强，$t$ 为实际气体温度，根据相似律写出 $Q$、$\eta$、$p$、$N$ 与 $Q_0$、$\eta_0$、$p_0$、$N_0$ 的换算关系式。
\end{xiaowen}
\xstarfull{1}\yiju{超纲：流体机械（泵与风机），赵琴第 11--14 章，复试范围；E5 题库问答题第 2 题}\nandu{超纲}

\end{ti}

\section*{三、计算题（4 题，共 58 分）}

\begin{ti}

\item （15 分）一弧形闸门，宽 $B=2\ \mathrm{m}$，圆心角 $\alpha=30^{\circ}$，半径 $r=3\ \mathrm{m}$，
闸门转轴 $O$ 与水面齐平。试求作用于闸门上的静水总压力 $F_x$、$F_z$，以及对转轴 $O$ 的总合力矩 $M$。
\tuyi{转轴 $O$ 位于水面上，同时是闸门圆弧的圆心；闸门为半径 $r=3\ \mathrm{m}$ 的圆柱面，
自水面附近的 $A$ 点向下张开圆心角 $\alpha=30^{\circ}$ 到 $B$ 点；闸门宽度（垂直于纸面）$B=2\ \mathrm{m}$；
上游水深恰为闸门的竖直投影高度。}
\xstarfull{5}\yiju{E4 2015 年计算"圆柱闸门静水总压力"；E5 题库计算题第 1 题（弧形闸门 $F_x$、$F_z$、$M$）；E1 slide33"第二章流体静力学＝计算题基础入门、重点掌握方法"}\nandu{中}

\item （15 分）水箱中的水通过垂直管道向大气出流。设水箱水深为 $H$，管道直径 $d$，长度 $l$，
沿程阻力系数 $\lambda$，局部阻力系数 $\zeta$。试求：① 在什么条件下流量 $Q$ 不随管长 $l$ 而变？
② 什么条件下流量 $Q$ 随管长 $l$ 的加大而增加？③ 什么条件下流量 $Q$ 随管长 $l$ 的加大而减小？
\xstarfull{4}\yiju{E4 2015 年计算题（突然缩小水头损失，水头损失计算）；E4 2026 年回忆版计算题第 5 题（管嘴体积流量，教材例题）；E1 slide33"第四章、第五章＝考试计算题重点"}\nandu{难}

\item （12 分）已知煤粉炉炉膛中上升烟气流的最小速度为 $v=0.5\ \mathrm{m/s}$，烟气的运动黏滞系数
$\nu=230\times10^{-6}\ \mathrm{m^2/s}$，问直径 $d=0.1\ \mathrm{mm}$ 的煤粉颗粒是沉降下来还是被烟气带走？
已知烟气的密度 $\rho=0.2\ \mathrm{kg/m^3}$，煤粉的密度 $\rho_m=1.3\times10^{3}\ \mathrm{kg/m^3}$。
\xstarfull{3}\yiju{E5 题库计算题第 3 题；E6 教材第 1 章牛顿内摩擦定律与第 5 章颗粒沉降（斯托克斯公式）；E4 历年真题未直接出现，故按题库型压至三星}\nandu{中}

\item （16 分）设有一台水泵，当转速 $n=1450\ \mathrm{r/min}$ 时，其性能参数列于下表：

\begin{center}
\begin{tabular}{@{}lccccccccc@{}}
\toprule
$Q$（$\mathrm{L/s}$） & 0 & 2 & 4 & 6 & 8 & 10 & 12 & 14 \\
\midrule
$H$（$\mathrm{m}$） & 11 & 10.8 & 10.5 & 10 & 9.2 & 8.4 & 7.4 & 6 \\
$\eta$（\%） & 0 & 15 & 30 & 45 & 60 & 65 & 55 & 30 \\
\bottomrule
\end{tabular}
\end{center}

管路系统的综合阻力数 $S=0.024\times10^{3}\ \mathrm{s^2/m^5}$，几何扬水高度 $H_Z=6\ \mathrm{m}$，
上下两水池水面均为大气压。
\begin{xiaowen}
\item 确定水泵运行时的工作参数（$Q$、$H$、$N$、$\eta$）。
\item 可以采取哪些方式调节运行工况？
\item 如以节流阀调节流量，使 $Q=8\ \mathrm{L/s}$ 时，有关工作参数是多少？
\end{xiaowen}
\xstarfull{2}\yiju{超纲：流体机械（泵与风机），赵琴第 11--14 章，复试范围；E5 题库计算题第 4 题；E1 slide33 对第 1--8 章的重点定位均未含本章}\nandu{难}

\end{ti}
